\begin{theorem}[12.2]\label{12.2}
    $E(\C)$~--- банахово, $A \in K(E), \lambda \in \C, \lambda \neq 0$. Тогда $\dim \ker A_\lambda < \infty$
\end{theorem}

% готовися к госу по диффурам
% я компакт
\begin{proof}
    Покажем, что единичная сфера в $\ker A_\lambda$~--- предкомпакт. Посмотрим на языке последовательностей
    \[ 
    \forall \{x_n \} \subset \ker A_\lambda:
    \begin{cases}
    \|x_n\| = 1, \text{ (на сфере) } \forall n\\
    A_{\lambda} x_n = 0
    \end{cases}
    \]
    Покажем, что из $\{x_n\}$ можно выделить сходящуюся подпоследовательность.
    $\{A x_n\}$~--- предкомпакт (так как $A$~--- компактный оператор), следовательно,
    по теореме 3.2 существует сходящаяся подпоследовательность $A x_{n_k} = \lambda x_{n_k} \to y$.

    \[\lambda x_{n_k} \to y \lra x_{n_k} \to \frac{1}{\lambda} y \implies Ax_{n_k} \to \frac{1}{\lambda} Ay\]

    Откуда получаем, что $Ay = \lambda y$ (единственность предела), а это значит, что $y \in \ker A_\lambda$.
\end{proof}

\begin{theorem}[12.3]\label{12.3}
$E(\C)$~--- банахово пространство, $A$~--- компактный оператор на $E$. Тогда
$\forall \delta > 0$ вне круга $\{ |\lambda| \leqslant \delta \}$ может лежать лишь конечное число собственных значений оператора $A$.

% тут должна быть картинка
\end{theorem}

\begin{example}
Пусть оператор в $l_2$ задаётся матрицей
$
\begin{pmatrix}
1 \\
_ & 0 \\
_ & _ & \frac{1}{2} \\
_ & _ & _ & 0 \\
_ & _ & _ & _ & \frac{1}{3} \\
_ & _ & _ & _ & _ & \ddots
\end{pmatrix}
$

Это компактный оператор, у которого $0$~--- собственное значение бесконечной кратности.
\end{example}

\begin{proof}
    Доказательство проведем для случая, когда $E$~--- гильбертово пространство, а $A$~--- компактный самосопряженный оператор. 

    Предположим противное. $\exists \delta_0$, что вне соответствующего круга $\{|\lambda| < \delta_0\}$ находится бесконечно много собственных значений оператора $A$. Выберем последовательность $\lambda_n$, из этого бесконечного множества. 
    Каждому $\lambda_n$ сопоставлен хотя бы один собственный вектор $e_n$. Тогда из самосопряженности оператора следует, что при $\lambda_n \neq \lambda_m$ выполнено $(e_n, e_m) = 0$. Отнормируем все $e_n$. 
    Следовательно, $A e_n$~--- предкомпактное множество (это эквивалентно вполне ограниченности).
    Покажем, что это не так. В самом деле, 
    \[\|Ae_n - Ae_m\|^2 = |\lambda_n|^2 + |\lambda_m|^2 > 2 \delta_0^2\]
    Отсюда следует, что $\{A e_n\}$~--- не вполне ограничено. Противоречие.
\end{proof}

\begin{lemmanote}
    В общем случае можно несложно доказать через теорему о почти перпендикуляре, но сейчас просто знакомимся с полезными приемами.
\end{lemmanote}
% обсуждаем шведов

\begin{corollary}[из теорем 12.2 и 12.3]
    \begin{enumerate}
        \item 
        \(
        \sum \limits_{|\lambda| > \delta > 0} \dim \ker A_{\lambda} < \infty
        \) (конечная сумма конечных размерностей)

        \item $\sigma_p(A)$~--- не более, чем счетное множество (нумеруем круги, вне каждого конечное множества, т.е. можем занумеровать мир ненулевых собственных значений)
    \end{enumerate}
\end{corollary}

\begin{theorem}[12.4 (Фредгольма)]
    Пусть $H(\C)$~--- гильбертово пространство, $A$~--- компактный самосопряжённый оператор, $\lambda \in \C, \lambda \neq 0$. Тогда $H = \im A_\lambda \oplus \ker A_\lambda$.
\end{theorem}
% Эта теорема -- как озеро Байкал, в нее втекает много маленьких теорем

% обсуждаем учебники по матанализу (Колмогоров vs Зорич)
\begin{lemma}[1]\label{fredholm:lemma1}
    Пусть $H(\C)$~--- гильбертово пространство, $A$~--- компактный самосопряжённый оператор, $\lambda \in \C, \lambda \neq 0$.
    Если $\lambda \in \sigma(A)$, то $\lambda$~--- собственное значение.
\end{lemma}

\begin{proof}
    По теореме \nameref{11.3} \[
    \lambda \in \sigma(A) \lra \exists \{x_n\} :
    \begin{cases}
        \|x_n\| = 1\ \forall n\\
        A_{\lambda} x_n \to 0
    \end{cases}
    \]
    Так как $A$~--- компактный, то $\{Ax_n\}$~--- предкомпакт (в нашем случае эквивалентно вполне ограниченности). Значит, существует сходящаяся подпоследовательность $\{Ax_{n_k}\}$, $Ax_{n_k} \to y$.
    Из второго условия получаем, что $A x_{n_k} - \lambda x_{n_k} \to 0$. Откуда $\lambda x_{n_k} \to y$. Получается $x_{n_k} \to \frac{1}{\lambda} y$, $A$~--- непрерывный $\implies Ax_{n_k} \to \frac{1}{\lambda} Ay$. Из единственности предела получаем, что $y = \frac{1}{\lambda} A y$, то есть $y \neq 0$~--- собственный вектор, а $\lambda$~--- собственное значение. Неравенство нулю следует из того, что мы все получали из сферы.
    
\end{proof}

\begin{lemma}[2]\label{fredholm:lemma2}
    Пусть $H(\C)$~--- гильбертово пространство, $A$~--- компактный самосопряжённый оператор, $\lambda \in \C, \lambda \neq 0$. Тогда $\overline{\im A_\lambda} = \im A_\lambda$
\end{lemma}


% Кирилл -- летописец
\begin{lemma}[3]\label{fredholm:lemma3}
    $A$~--- самосопряжённый оператор в $H(\C)$, $M$ (подпространство) инвариатно относительно $A$, то есть $A M \subset M$. Тогда $M^\perp$ инвариантен относительно $A$.
\end{lemma}

\begin{proof}
    $y \in M^\perp$. Возьмем произвольный $x \in M$ и рассмотрим $(x, Ay) = (Ax, y)$. Так как $M$ инвариантно, то $Ax \in M$, а значит $(x, Ay)=(Ax, y)=0$, то есть $Ay \in M^\perp$ и $M^\perp$ инвариантно относительно $A$
\end{proof}

\begin{proof}[леммы 2]
    \begin{enumerate}
        \item Ясно, что $\ker A_{\lambda}$ инвариантно относительно $A$ (так как образовано собственными векторами) %~--- говорит Андрей.
        Отсюда по лемме \nameref{fredholm:lemma3} следует инвариантность $\overline{\im A_\lambda}$, т.к. из теоремы \nameref{11.2} мы знаем, что $\overline{\im A_{\lambda}} \oplus_{\perp} \ker A_{\lambda} = H$.
        \item Рассмотрим $\tilde{A}=A|_{\overline{\im A_\lambda}}$~--- компактный самосопряженный оператор. 
        
        $\tilde{A}_\lambda x = y, x \in \overline{\im A_\lambda} \implies y \in \im \tilde{A_\lambda} \subset \im A_\lambda$
        
        У него нет собственных векторов (они все остались во втором слагаемом прямой суммы: если это собственное значение, то соответствующий собственный ветор $\tilde{A}$ лежащий в $\overline{\im A_\lambda}$. Тогда он бы являлся и собственным вектором $A$, но мы их отбросили - противоречие), то есть число $\lambda$ не является собственным значением оператора $\tilde{A}$. Отсюда, по лемме \nameref{fredholm:lemma1}, $\lambda \not\in \sigma(\tilde{A})$. Следовательно, $\lambda \in \rho(\tilde{A})$ и $\tilde{A_\lambda}x=y$ имеет решение $\implies \exists x \in \overline{\im A_\lambda} \subset H: \tilde{A_\lambda}x=y$.
        
        Получаем, что $\overline{\im \tilde{A_{\lambda}}} = \im \tilde{A_{\lambda}}$.
        Осталось заметить, что действия операторов $A$ и $\tilde{A}$ на множестве $\im A$ совпадают, а значит последнее равенство верно и для оператора $A_{\lambda}$.
    \end{enumerate}
\end{proof}

\begin{proof}[теоремы Фредгольма]
    По \nameref{11.2}: $H=\overline{\im A_\lambda} \oplus \ker A_\lambda$. По лемме \nameref{fredholm:lemma2}: $\overline{\im A_\lambda} = \im A_\lambda$.
\end{proof}

% карту ручейков нарисуем в следующий раз